\documentclass[10pt,a4paper]{amsart}
\usepackage[margin=19mm]{geometry}
\usepackage{amsmath,amssymb,mathtools}
\usepackage[scr=rsfs,cal=euler]{mathalfa}
\usepackage{xfrac}
\usepackage[unicode,hidelinks,bookmarks=false]{hyperref}
\newcommand{\ie}{i.e.\ }
\newcommand{\cf}{cf.\ }
\newcommand{\resp}{respectively\ }
\newcommand{\Subsection}{\subsection}
\providecommand{\nobreakdash}{\nobreak\hbox{-}}
\setlength{\emergencystretch}{2em}
\multlinegap0pt
\usepackage{bm}
\usepackage{bbm}
\usepackage{dsfont}
\usepackage{xspace}
\usepackage{cancel}
\usepackage{mathtools}
\usepackage{amsthm}
\makeatletter
\@ifundefined{assumption}{\newtheorem{assumption}{Assumption}}{}
\@ifundefined{lemma}{\newtheorem{lemma}{Lemma}}{}
\makeatother
\newcommand{\E}{\mathbb{E}}
\newcommand{\R}{\mathbb{R}}
\newcommand{\bU}{\mathbf{U}}
\newcommand{\cH}{\mathcal{H}}
\newcommand{\dd}{\mathrm{d}}
\title[A Zeroth-Order Deep Learning Method for Fully Nonlinear Para]{A Zeroth-Order Deep Learning Method for Fully Nonlinear Parabolic Partial Differential Equations with Unknown Coefficients}
\author{Yanwei Jia, Du Ouyang, Huy\^en Pham, Xun Yu Zhou}
\date{Source: arXiv:2606.24999v1 (CC BY 4.0)}
\begin{document}
\maketitle

\noindent\emph{Source and licence.} A Zeroth-Order Deep Learning Method for Fully Nonlinear Parabolic Partial Differential Equations with Unknown Coefficients. Version arXiv:2606.24999v1 (CC BY 4.0), distributed under the Creative Commons Attribution 4.0 International licence, as recorded in the arXiv licence metadata for this exact version and shown on the abstract page https://arxiv.org/abs/2606.24999v1. Full rights notice: licenses/arxiv-2606.24999-CC-BY-4.0.txt.\par\smallskip

\noindent\textit{Editorial scope.} Counted unit: the Lemma whose source label is \texttt{lemma: discretization error} in the article (source0.tex lines 2230--2235, source label \texttt{lemma: discretization error}), delivered together with the article's own complete original proof of it (source0.tex lines 2236--2393, one \texttt{proof} environment of 158 lines). No mathematics is written, repaired or re-derived: statement and proof are the article's own text line for line. Required context retained uncounted: eq:sde (lines 260--262); assum: sde coef (lines 500--507); assum: hypothesis space (lines 819--822); assum: reward and f regularity (lines 834--838). Every definition, display and label of those lines is retained unchanged. Omitted from this package and not redistributed: 1--259, 263--499, 508--818, 823--833, 839--2229, 2394--3076. Nothing inside a retained excerpt is omitted or altered: every retained body is delivered exactly as the article's lines are written, so no omission receipt of any kind accompanies this package. Citations: the retained excerpts carry no citation command at all, so no bibliography entry of the article is transcribed and no reference is redistributed. Reference adaptation: none was needed and none was made; every reference inside the delivered unit resolves inside it, so no unresolved reference or citation is produced. Printed slips kept literal: no printed word, symbol, hypothesis or number of the counted statement or of its proof was altered. Layout and class: the article's own class is \texttt{article} with packages including jmlr2e, bbm, hyperref, amsmath, amssymb, mathrsfs, caption, subcaption, booktabs, multicol, multirow, xcolor, color, algorithm; the wrapper is the lane's pinned \texttt{amsart} editorial wrapper at 10pt on A4, which restates the theorem-like environments the retained text needs and adds the article's own macro definitions that the retained text needs (\texttt{\textbackslash E}, \texttt{\textbackslash R}, \texttt{\textbackslash bU}, \texttt{\textbackslash cH}, \texttt{\textbackslash dd}), with extra wrapper packages (bm, bbm, dsfont, xspace, cancel, mathtools). The delivered numbering is the unit's own. Assets: no retained span contains a figure environment, a graphics-inclusion command, a PDF-page inclusion, a table or a TikZ picture; no image file is copied into this package, the package declares no asset dependency and no image file is redistributed with it. Licence: version arXiv:2606.24999v1 of this article is distributed under the Creative Commons Attribution 4.0 International licence, as recorded in the arXiv licence metadata for this exact version and shown on the abstract page https://arxiv.org/abs/2606.24999v1. A full rights notice is delivered at licenses/arxiv-2606.24999-CC-BY-4.0.txt. Characters: every retained excerpt is pure ASCII, so the delivered engine, XeLaTeX with its default Latin Modern text font, reports no missing character.
\begin{equation}\label{eq:sde}
     \dd X_s = b(X_s) \dd s + \sqrt{2} \sigma(X_s) \dd W_s,
\end{equation}

\begin{assumption}\label{assum: sde coef}
The coefficients $b$ and $\sigma$ are sufficiently regular so that the SDE \eqref{eq:sde} admits a unique strong solution. Moreover, they satisfy the linear growth and bounded condition, i.e., there are constants $ L_b$ and $ \sigma_0$ such that
    \begin{equation}
        \begin{aligned}
            |b(t,x)| \le L_b(1+|x|), \quad |\sigma(t,x)| \le \sigma_0, \quad \text{for all}\ \ (t,x)\in [0,T]\times\R^d.
        \end{aligned}
    \end{equation}
\end{assumption}

\begin{assumption}
\label{assum: hypothesis space}
The activation function satisfies $\rho\in C_b^5(\mathbb R)$.  The parameter set $\Theta\subset\mathbb R^p$ is compact and satisfies $|\theta|_\infty\le R$ for all $\theta\in\Theta$, where $p$ is the total number of trainable parameters in the three networks.
\end{assumption}

\begin{assumption}
\label{assum: reward and f regularity}
The functions $J_1(\cdot, \cdot; \varphi)\in C_p^{1,4}([0,T]\times\R^d), J_2(\cdot,\cdot;\varphi)\in C_p^{1,2}([0,T]\times \R^d), \tilde f \in C_p^{1,2}([0,T]\times \R^d)$ uniformly over the hypothesis space $ \varphi\in \mathcal{H}$ in the sense that their derivatives have polynomial growth with constants
independent of $\varphi$. The terminal function $g$ has polynomial growth.
\end{assumption}

\begin{lemma}\label{lemma: discretization error}
Under Assumptions \ref{assum: hypothesis space}, \ref{assum: reward and f regularity}, \ref{assum: sde coef}, there is a constant $ C$ independent of $ \epsilon$ and $ N$ such that
\begin{equation}
    \sup_{\varphi \in \cH} \left|  L^{\mathscr{G}}(\varphi) - L(\varphi)\right| \le C \epsilon^{-4}N^{-1}.
\end{equation}
\end{lemma}

\begin{proof}
    We aim to bound the difference between the continuous-time population loss $L(\varphi)$ and its discrete-time approximation $L^{\mathcal{G}}(\varphi)$. We prove the Hessian loss for illustration, and the value and gradient components are treated in the same way and are of lower order in $\epsilon$. Let the estimator be denoted by
\begin{equation}
    \xi^\epsilon(Z) := \frac{ZZ^\top - I_d}{\epsilon^2},\quad \tilde f(t,x):=f(t,x,\mathbf U_{n-1}(t,x)).
\end{equation}
The continuous and discrete losses are defined as:
\begin{align}
    L_h(\varphi_h) &:=  \int_0^T \E \left[ e^{\beta t} \left| \varphi_h(t, X_t) - \xi^\epsilon(Z)R^{t, X_t+\epsilon Z}(\mathbf{U}_{n-1}) \right|^2 \right] \dd t, \\
    L_h^{\mathcal{G}}(\varphi_h) &:=  \sum_{i=0}^{N-1} \E \left[ e^{\beta t_i} \left| \varphi_h(t_i, X_{t_i}) - \xi^\epsilon(Z) \tilde R^{t_i, X_{t_i}+\epsilon Z}(\mathbf{U}_{n-1}) \right|^2 \right] \Delta t_i,
\end{align}
where $\tilde R^{t,x}(\bU_{n-1})$ uses the discrete reward approximation
\begin{equation}
    \tilde R^{t,x}(\bU_{n-1}) = g(X_T^{t,x}) + \sum_i f(t_i,X_{t_i}^{t,x},\bU_{n-1}(t_i,X_{t_i}^{t,x}))\Delta t_i,
\end{equation}
instead of the exact $R^{t,x}(\bU_{n-1})$
\begin{equation}
    R^{t,x}(\bU_{n-1}) = g(X_T^{t,x}) + \int_t^{T} f(s,X_s^{t,x},\bU_{n-1}(s,X_{s}^{t,x}))\dd s.
\end{equation}
Denote
\[
    \Delta(t,y)
    :=
    \tilde R^{t,y}(\mathbf U_{n-1})
    -
    R^{t,y}(\mathbf U_{n-1}).
\]
Introduce the intermediate grid loss $\bar L_h^{\mathscr G}$ obtained by
replacing the time integral in $L_h$ with the grid sum while keeping the exact
reward $R$.
\begin{equation}
    \Bar{L}_h(\varphi_h)^{\mathscr{G}} :=   \sum_{i=0}^{N-1} \E \left[ e^{\beta t_i} \left| \varphi_h(t_i, X_{t_i}) - \xi^\epsilon(Z)  R^{t_i, X_{t_i}+\epsilon Z}(\mathbf{U}_{n-1}) \right|^2 \right] \Delta t_i.
\end{equation}

We have the following inner and outer error decompositions
\begin{equation}
    \begin{aligned}
        \left|L_h(\varphi_h) - L_h^{\mathcal{G}}(\varphi_h)\right|\le \underbrace{\left|L_h^{\mathscr G}(\varphi_h)-\bar L_h^{\mathscr G}(\varphi_h) \right|}_{\textrm{Inner Discretization Error}}+ \underbrace{\left| \bar L_h^{\mathscr G}(\varphi_h)-L_h(\varphi_h)\right|}_{\textrm{Outer Discretization Error}}.
    \end{aligned}
\end{equation}
Therefore, there are two types of error -- the inner and the outer discretization errors.
For the inner discretization error, it suffices to bound
\begin{equation}\label{eq: inner decomp}
\begin{aligned}
     &\left| \E \left[ \left\| \varphi_h(t_i, X_{t_i}) - \xi^\epsilon(Z)R^{t_i, X_{t_i}+\epsilon Z} \right\|^2 \right] - \E \left[ \left\| \varphi_h(t_i, X_{t_i}) - \xi^\epsilon(Z)\tilde R^{t_i, X_{t_i}+\epsilon Z} \right\|^2 \right]\right|\\
     &\le \E\left[ \left| 2\varphi_h(t_i,X_{t_i}) - \xi^{\epsilon}(Z)(\tilde R^{t_i,X_{t_i}+ \epsilon Z} + R^{t_i,X_{t_i} + \epsilon Z})\right|\left| \xi^{\epsilon}(Z)\right|\left|\Delta(t_i,X_{t_i}+\epsilon Z)\right|\right].
\end{aligned}
\end{equation}
Due to the polynomial growth condition and the moment bounds of $ X_t$, it suffices to bound $ \|\Delta(t_i,X_{t_i} + \epsilon Z)\|_{L_2}$.
To this end,
recall that $R^{t,y} = g(X_T^{t,y}) + \int_t^T \tilde f(s, X_s^{t,y}) \dd s$. The terminal part cancels and
\[
    \Delta(t_i,y)
    =
    \sum_{k=i}^{N-1}
    \int_{t_k}^{t_{k+1}}
    \left[
    \tilde f(t_k,X_{t_k}^{t_i,y})
    -
    \tilde f(s,X_s^{t_i,y})
    \right]\dd s .
\]
By Assumption \ref{assum: reward and f regularity}, It\^o's formula gives
\[
    \dd \tilde f(s,X_s^{t_i,y})
    =
    A_s\,\dd s+ B_{s}\,\dd W_s,
    \qquad
    \E\left[|A_s|^2+|B_{s}|^2\right]
    \le C(1+|y|^q)
\]
for some $q\ge0$, uniformly in $s$.
Stochastic Fubini therefore yields
\[
\begin{aligned}
    \Delta(t_i,y)
    =
    -\sum_{k=i}^{N-1}
    \int_{t_k}^{t_{k+1}}(t_{k+1}-u)A_u\,\dd u
    -\sum_{k=i}^{N-1}
    \int_{t_k}^{t_{k+1}}(t_{k+1}-u)B_{u}\,\dd W_u .
\end{aligned}
\]
The drift term has $L^2$ norm at most $CN(T/N)^2\le C/N$.
For the martingale term, It\^o's isometry and orthogonality of martingale
increments give
\[
    \E\left|
    \sum_{k=i}^{N-1}
    \int_{t_k}^{t_{k+1}}(t_{k+1}-u)B_{u}\,\dd W_u
    \right|^2
    \le
    CNh^3
    \le
    \frac{C}{N^2}.
\]
Hence
\[
    \|\Delta(t_i,y)\|_{L^2}
    \le
    \frac{C(1+|y|^q)}{N}.
\]
Conditioning on $Z$, taking $y=X_{t_i}+\epsilon Z$, and using the moment
bounds for $X_{t_i}$ together with Gaussian moments,
\[
    \left\|
    \Delta(t_i,X_{t_i}+\epsilon Z)
    \right\|_{L^2}
    \le
    \frac{C}{N}.
\]
This along with \eqref{eq: inner decomp} yields
\[
    \sup_{\varphi_h\in\mathcal H}
    |L_h^{\mathscr G}(\varphi_h)-\bar L_h^{\mathscr G}(\varphi_h)|
    \le
    \frac{C}{N\epsilon^4}.
\]


Now we bound the outer discretization error:
\begin{equation}
    \left| \int_0^T l(t) \dd t - \sum_{i=0}^{N-1} l(t_i) \Delta t_i \right|, \quad \text{where } l(t) = e^{\beta t} \E \left[ \left\| \varphi_h(t, X_t) - \xi^\epsilon(Z)R^{t,X_t + \epsilon Z} \right\|^2 \right].
\end{equation}
It suffices to show that the time derivative $\frac{\dd}{\dd t} l(t)$ is bounded. Indeed,
\begin{equation}
\begin{aligned}
    l(t) &= e^{\beta t}\E \left[ \left\| \varphi_h(t, X_t) - \xi^\epsilon(Z)R^{t,X_t + \epsilon Z} \right\|^2 \right] \\
    &= e^{\beta t}\left(\E\left[ \left\|\varphi_h(t,X_t)\right\|^2\right] + \E \left[ \left\| \xi^{\epsilon}(Z) R^{t, X_t+\epsilon Z} \right\|^2 \right] - 2\E\left[ \langle \varphi_h(t,X_t), \xi^{\epsilon}(Z) \rangle R^{t,X_t + \epsilon Z}\right]\right).
    \end{aligned}
\end{equation}
Since $ \varphi_h$ is smooth and bounded, we need only to analyze the second and third terms.

For the second term, using the conditional expectation, we have
\begin{equation}
    I_2(t):=\E\left[ \left\| \xi^{\epsilon}(Z) R^{t, X_t+\epsilon Z} \right\|^2\right] = \E\left[ \|\xi^{\epsilon}(Z)\|^2 \E\left[ (R^{t,X_t + \epsilon Z})^2\mid Z,X_t\right]\right].
\end{equation}

Denote $u(t, x) := \E [ (R^{t,x})^2 ]$.
Under Assumption \ref{assum: reward and f regularity}, the  derivatives of $ u$ satisfy the polynomial growth condition.
Therefore, the time derivative of the expectation is:
\begin{equation}
\begin{aligned}
    \frac{\dd}{\dd t}\E\left[ \left\| \xi^{\epsilon}(Z) R^{t, X_t+\epsilon Z} \right\|^2\right]& =   \E[\|\xi^{\epsilon}(Z)\|^2 \frac{\dd}{\dd t}u(t, X_t + \epsilon Z)] \\
    &= \E \left[ \|\xi^{\epsilon}(Z)\|^2(\partial_t + \mathcal{L}) F_Z(t,x)\big|_{x = X_t} \right],
    \end{aligned}
\end{equation}
where $ F_Z(t,x) = u(t,x + \epsilon Z)$.
This confirms that $|\frac{\dd}{\dd t} I_2(t)| \le C \epsilon^{-4}$.

Similarly, for the cross term involving $\E [ \langle \varphi, \xi^\epsilon(Z)R^{t,X_t+\epsilon Z} \rangle ]$, the derivative involves terms related to $\E[R^{t,x}]$. Using Assumption \ref{assum: reward and f regularity}, the derivative is bounded by $C \epsilon^{-2}$.


Combining these bounds, the total discretization error is bounded by:
\begin{equation}
    \sup_{\varphi \in \cH} \left| L_h^{\mathcal{G}}(\varphi) - L_h(\varphi) \right| \le C \left( \epsilon^{-4} N^{-1} + \epsilon^{-2} N^{-1} \right) \le C \epsilon^{-4} N^{-1}.
\end{equation}
This proves the desired result.
\end{proof}

\end{document}
